\documentclass[11pt]{article}
\usepackage[margin=1in]{geometry}
\usepackage[T1]{fontenc}
\usepackage[utf8]{inputenc}
\usepackage[french]{babel}
\usepackage{amsmath,amssymb,amsthm}
\usepackage[colorlinks=true,linkcolor=blue,urlcolor=blue]{hyperref}
\newtheorem{problem}{Problème}
\newenvironment{refs}{\par\small\noindent\emph{Références :}\begin{itemize}\setlength\itemsep{0pt}}{\end{itemize}\normalsize}
\title{Hors de la Caverne \\ \Large Analyse}
\author{}
\date{}
\begin{document}
\maketitle
\noindent\emph{Des problèmes, pas de réponses.} Chaque problème ci-dessous est posé
sans solution. Les références indiquent où trouver les connaissances nécessaires.
\medskip


\section*{Lycée}

\begin{problem}[{La limite, à partir des premiers principes --- Lycée}]
\textbf{CAL-01.} \textbf{Problème.} Calculez
$ \lim_{x\to 1} \frac{x^2-1}{x-1} $,
en expliquant pourquoi la réponse n'est pas \og{} $0/0$ \fg{} et pourquoi simplifier par $ x-1 $ est
légitime bien que la fonction ne soit pas définie en $ x=1 $. Énoncez ensuite la définition
$\varepsilon$--$\delta$ de $ \lim_{x\to a} f(x)=L $, et démontrez directement à partir d'elle :
\begin{enumerate}
\item[(a)] que la limite que vous
avez calculée est correcte ;
\item[(b)] la loi de la somme
$ \lim (f+g) = \lim f + \lim g $ et la loi du produit pour les limites --- les deux lois que
tout calcul ultérieur de cette page utilise en silence.
\end{enumerate}

\end{problem}

\noindent Newton et Leibniz calculaient avec des quantités \og{} infiniment petites \fg{}
et furent moqués pour cela --- Berkeley (1734) appelait leurs accroissements évanouissants
\og{} les fantômes de quantités défuntes \fg{}. La réponse moderne a demandé 150 ans : le
\textit{Cours d'analyse} de Cauchy (1821) a placé la limite au fondement, et la formulation
$\varepsilon$--$\delta$ de Weierstrass (cours berlinois des années 1860) l'a rendue inattaquable. Démontrer les lois
des limites une fois pour toutes est le droit d'entrée à tout ce qui suit.

\begin{refs}
\item Contexte : Limite d'une fonction --- Wikipédia (\url{https://fr.wikipedia.org/wiki/Limite_%28math%C3%A9matiques_%C3%A9l%C3%A9mentaires%29})
\item Lecture : Michael Spivak, \textit{Calculus}, ch. 5 (les limites)
\item Cours : MIT OCW 18.01SC Single Variable Calculus (\url{https://ocw.mit.edu/courses/18-01sc-single-variable-calculus-fall-2010/})
\item Vidéo : 3Blue1Brown, Essence of Calculus (\url{https://www.3blue1brown.com/lessons/essence-of-calculus/})
\end{refs}

\bigskip

\begin{problem}[{La dérivée est une limite --- Lycée}]
\textbf{CAL-02.} \textbf{Problème.} Posez
\[ f'(a) = \lim_{h\to 0} \frac{f(a+h)-f(a)}{h}. \]
Travaillez uniquement à partir de cette définition --- aucune règle.
\begin{enumerate}
\item[(a)] Calculez les dérivées de $ f(x)=x^2 $, $ f(x)=\sqrt{x} $ et
$ f(x)=1/x $.
\item[(b)] Démontrez la règle du produit $ (fg)' = f'g + fg' $ directement à partir de la définition, en précisant
exactement quelles lois des limites de CAL-01 vous invoquez.
\item[(c)] Démontrez que la dérivabilité de $ f $ en un point y entraîne sa continuité ---
la partie (b) en a besoin.
\end{enumerate}

\end{problem}

\noindent Les premiers manuscrits privés de Leibniz (1675) ont brièvement supposé que
$ d(uv) = du \cdot dv $ avant qu'il ne se corrige --- la règle du produit est bel et bien un
théorème, non une trivialité de notation. Newton appelait les dérivées des \og{} fluxions \fg{} de
\og{} fluentes \fg{} ; la définition par le taux d'accroissement ci-dessus est celle de Cauchy. Faire
ces calculs honnêtement, une fois, est ce qui sépare le fait de savoir l'analyse du fait de s'en souvenir.

\begin{refs}
\item Contexte : Dérivée --- Wikipédia (\url{https://fr.wikipedia.org/wiki/D%C3%A9riv%C3%A9e})
\item Lecture : Michael Spivak, \textit{Calculus}, ch. 9--10
\item Cours : MIT OCW 18.01SC Single Variable Calculus (\url{https://ocw.mit.edu/courses/18-01sc-single-variable-calculus-fall-2010/})
\item Vidéo : 3Blue1Brown, Essence of Calculus (\url{https://www.3blue1brown.com/lessons/essence-of-calculus/})
\end{refs}

\bigskip

\begin{problem}[{Les tangentes avant l'analyse --- Lycée}]
\textbf{CAL-03.} \textbf{Problème.} Travaillez uniquement avec de l'algèbre --- ni limites, ni dérivées.
\begin{enumerate}
\item[(a)] Trouvez la tangente à la parabole $ y=x^2 $ au point $ (a, a^2) $ : exigez
que la droite $ y = mx + c $ rencontre la parabole en une racine double en $ x=a $, et résolvez
en $ m $ et $ c $.
\item[(b)] Faites de même pour $ y=x^3 $ en $ (a, a^3) $.
\item[(c)] Vérifiez que les deux réponses s'accordent avec la dérivée, puis expliquez précisément pourquoi
la méthode de la racine double fonctionne pour les polynômes mais ne peut pas définir la tangente à
$ y=\sin x $.
\end{enumerate}

\end{problem}

\noindent C'est ainsi que l'on trouvait les tangentes avant l'existence de l'analyse. La
méthode des normales de Descartes (\textit{La Géométrie}, 1637) et l'\textit{adégalité} de Fermat (diffusée
dans les années 1630) résolvaient des problèmes de tangente et d'extremum pour les courbes algébriques une génération avant que
Newton ne naisse ; leur limitation aux courbes algébriques est exactement ce que la notion de limite
a levé. La méthode de Fermat anticipe la dérivée d'assez près pour que Lagrange lui en ait attribué
l'invention.

\begin{refs}
\item Contexte : Adégalité --- Wikipédia (\url{https://fr.wikipedia.org/wiki/Ad%C3%A9galit%C3%A9})
\item Contexte : Tangente --- Wikipédia (\url{https://fr.wikipedia.org/wiki/Tangente_%28g%C3%A9om%C3%A9trie%29})
\item Lecture : Carl B. Boyer, \textit{The History of the Calculus and Its Conceptual Development}
\end{refs}

\bigskip

\begin{problem}[{Le plus grand champ qu'une clôture puisse enfermer --- Lycée}]
\textbf{CAL-04.} \textbf{Problème.} Vous disposez de $ L $ mètres de clôture.
\begin{enumerate}
\item[(a)] Démontrez que
parmi tous les rectangles de périmètre $ L $, le carré enferme la plus grande aire --- deux fois :
une fois avec l'analyse, et une fois sans analyse du tout, à l'aide de l'inégalité arithmético-géométrique
$ \sqrt{xy} \le \frac{x+y}{2} $ (démontrez-la également).
\item[(b)] Cette fois la clôture ferme un
champ rectangulaire contre le long mur droit d'une grange, de sorte que seuls trois côtés sont à clôturer :
trouvez les proportions optimales, avec démonstration.
\item[(c)] Expliquez pourquoi trouver un point où
$ f'=0 $ n'est \textit{pas} à soi seul une démonstration d'optimalité, et complétez la justification
dans les deux cas (comportement au bord, dérivée seconde, ou concavité).
\end{enumerate}

\end{problem}

\noindent L'optimisation est la plus ancienne publicité pour l'analyse. L'intuition
isopérimétrique --- plus c'est symétrique, plus cela contient --- remonte à l'Antiquité (le problème de Didon dans
l'\textit{Énéide} de Virgile ; Zénodore, vers 180 av. J.-C.), et la \textit{Nova stereometria
doliorum vinariorum} de Kepler (1615) optimisait les proportions des tonneaux avant l'existence des dérivées.
La partie (iii) est l'âme du problème : un point critique est un suspect, pas un verdict.

\begin{refs}
\item Contexte : Théorème isopérimétrique --- Wikipédia (\url{https://fr.wikipedia.org/wiki/Th%C3%A9or%C3%A8me_isop%C3%A9rim%C3%A9trique})
\item Lecture : James Stewart, \textit{Calculus}, ch. 4 (applications de la dérivation)
\item Lecture : Michael Spivak, \textit{Calculus}, ch. 11
\end{refs}

\bigskip

\begin{problem}[{La boîte de conserve la plus économique --- Lycée}]
\textbf{CAL-05.} \textbf{Problème.} Une boîte cylindrique fermée doit contenir un volume fixé
$ V $. Démontrez que l'aire totale de sa surface (le côté plus les deux couvercles) est minimale exactement lorsque
la hauteur égale le diamètre, $ h = 2r $. Justifiez que votre point critique est un véritable
minimum global sur $ (0,\infty) $ --- examinez ce qu'il advient de l'aire lorsque $ r\to 0^+ $
puis lorsque $ r\to\infty $. Mesurez ensuite une vraie boîte de votre cuisine et déterminez à quelle distance
de l'optimum elle se trouve ; suggérez, en une phrase, quels coûts autres que la matière pourraient expliquer
l'écart.
\end{problem}

\noindent Le problème de la boîte est l'optimisation sous contrainte classique à une variable :
une contrainte (le volume) élimine une variable et laisse une honnête minimisation à une
variable. Les vraies boîtes sont nettement plus hautes que l'optimum mathématique --- une leçon précoce :
un modèle démontre des théorèmes sur lui-même, et l'étape de modélisation est une responsabilité
distincte.

\begin{refs}
\item Contexte : Optimisation (mathématiques) --- Wikipédia (\url{https://fr.wikipedia.org/wiki/Optimisation_%28math%C3%A9matiques%29})
\item Lecture : James Stewart, \textit{Calculus}, ch. 4.7 (problèmes d'optimisation)
\item Cours : MIT OCW 18.01SC Single Variable Calculus (\url{https://ocw.mit.edu/courses/18-01sc-single-variable-calculus-fall-2010/})
\end{refs}

\bigskip

\begin{problem}[{Les taux liés, honnêtement --- Lycée}]
\textbf{CAL-06.} \textbf{Problème.} Une échelle de 5~m est appuyée contre un mur ; son pied
glisse en s'éloignant à 1~m/s.
\begin{enumerate}
\item[(a)] Trouvez la vitesse du sommet lorsque le pied est à 3~m
du mur, en précisant exactement l'étape de dérivation des fonctions composées que votre calcul utilise et le théorème qui
l'autorise.
\item[(b)] Montrez que lorsque le pied s'approche de 5~m du mur, votre formule
prédit que la vitesse du sommet tend vers l'infini ; expliquez ce que cela démontre au sujet du modèle
(le sommet d'une vraie échelle reste-t-il contre le mur ?).
\item[(c)] De l'eau s'écoule d'un
cône renversé de demi-angle 30$^\circ$ : reliez $ dV/dt $ au taux de variation de la hauteur d'eau,
en dérivant --- et non en citant --- la relation de la formule du volume $ V = \tfrac{1}{3}\pi r^2 h $
que vous utilisez.
\end{enumerate}

\end{problem}

\noindent Les problèmes de taux liés sont un classique depuis les manuels britanniques du
dix-neuvième siècle issus de la tradition des \og{} mathématiques mixtes \fg{}. Le paradoxe de l'échelle
qui tombe, en (ii), est un échec véritablement instructif : la mathématique est correcte, c'est la
prémisse (le sommet collé au mur) qui cède --- et l'analyse elle-même détecte le moment
où elle doit céder.

\begin{refs}
\item Contexte : Taux liés --- Wikipédia (en anglais) (\url{https://en.wikipedia.org/wiki/Related_rates})
\item Contexte : Dérivation des fonctions composées --- Wikipédia (\url{https://fr.wikipedia.org/wiki/Th%C3%A9or%C3%A8me_de_d%C3%A9rivation_des_fonctions_compos%C3%A9es})
\item Lecture : James Stewart, \textit{Calculus}, ch. 3.9 (taux liés)
\end{refs}

\bigskip

\begin{problem}[{Le théorème du compteur de vitesse --- Lycée}]
\textbf{CAL-07.} \textbf{Problème.} Une voiture passe un premier péage à midi et un second,
90~km plus loin, à 13~h. Démontrez qu'à un certain instant son compteur indiquait
exactement 90~km/h. Pour le faire correctement : énoncez le théorème de Rolle ; déduisez-en le
théorème des accroissements finis --- pour $ f $ convenable il existe $ c\in(a,b) $ tel que
$ f'(c) = \frac{f(b)-f(a)}{b-a} $ --- et énoncez les hypothèses de modélisation sur la fonction
position dont votre démonstration a besoin. Utilisez ensuite ce théorème pour démontrer qu'une fonction dont la dérivée
s'annule sur un intervalle y est constante, et expliquez en un paragraphe pourquoi toute la
théorie des primitives (le \og{} $ +\,C $ \fg{}) repose sur ce fait d'apparence anodine.
\end{problem}

\noindent Michel Rolle a énoncé son théorème pour les polynômes en 1691 --- tout en étant
l'un des critiques les plus acérés de l'analyse ; il la traitait de recueil de sophismes ingénieux.
Cauchy a fait du théorème des accroissements finis le moteur de l'analyse rigoureuse : c'est le pont qui mène de
l'information locale (la dérivée en chaque point) aux conclusions globales (ce que la fonction
a fait sur tout l'intervalle). Presque tous les théorèmes de cette bande traversent ce pont.

\begin{refs}
\item Contexte : Théorème des accroissements finis --- Wikipédia (\url{https://fr.wikipedia.org/wiki/Th%C3%A9or%C3%A8me_des_accroissements_finis})
\item Contexte : Théorème de Rolle --- Wikipédia (\url{https://fr.wikipedia.org/wiki/Th%C3%A9or%C3%A8me_de_Rolle})
\item Lecture : Michael Spivak, \textit{Calculus}, ch. 11
\item Cours : MIT OCW 18.01SC Single Variable Calculus (\url{https://ocw.mit.edu/courses/18-01sc-single-variable-calculus-fall-2010/})
\end{refs}

\bigskip

\begin{problem}[{L'aire à partir des premiers principes --- Lycée}]
\textbf{CAL-08.} \textbf{Problème.} En utilisant la définition de l'intégrale comme limite de
sommes de Riemann, calculez $ \int_0^1 x^2\,dx $ directement : découpez $ [0,1] $ en $ n $
morceaux égaux, écrivez les sommes supérieures et inférieures, évaluez-les à l'aide de l'identité
\[ 1^2 + 2^2 + \cdots + n^2 = \frac{n(n+1)(2n+1)}{6} \]
(démontrez-la par récurrence), et montrez que les deux sommes convergent vers la même limite. Calculez ensuite
$ \int_0^b x^k\,dx $ pour tout entier $ k \ge 1 $ par l'astuce de Fermat : découpez
$ [0,b] $ non pas régulièrement mais aux points géométriques $ b, br, br^2, \dots $ avec
$ r<1 $, sommez les rectangles obtenus comme une série géométrique, et faites tendre $ r\to 1 $.
\end{problem}

\noindent Archimède a obtenu l'aire sous une parabole vers 250 av. J.-C. par la
méthode d'exhaustion --- la même réponse, deux millénaires avant que les limites ne soient définies. L'astuce
du découpage géométrique de Fermat (années 1630) a quadraturé toutes les courbes $ y=x^k $ d'un coup, encore
sans aucun théorème fondamental : l'intégration est plus ancienne que la dérivation, et la pratiquer
une fois à la dure est ce qui vous fait sentir ce que le théorème fondamental (CAL-09) épargne réellement.

\begin{refs}
\item Contexte : Somme de Riemann --- Wikipédia (\url{https://fr.wikipedia.org/wiki/Somme_de_Riemann})
\item Contexte : La Quadrature de la parabole --- Wikipédia (\url{https://fr.wikipedia.org/wiki/La_Quadrature_de_la_parabole_%28Archim%C3%A8de%29})
\item Lecture : Tom M. Apostol, \textit{Calculus}, vol. 1, ch. 1
\end{refs}

\bigskip

\begin{problem}[{Pourquoi le théorème fondamental est un théorème --- Lycée}]
\textbf{CAL-09.} \textbf{Problème.} Énoncez précisément les deux parties du théorème fondamental
de l'analyse, avec leurs hypothèses.
\begin{enumerate}
\item[(a)] Démontrez la partie 1 : si $ f $ est continue sur
$ [a,b] $ et $ F(x) = \int_a^x f(t)\,dt $, alors $ F'(x) = f(x) $ --- servez-vous de la continuité
de $ f $ et encadrez le taux d'accroissement entre le minimum et le maximum de $ f $
sur un intervalle qui rétrécit.
\item[(b)] Déduisez-en la partie 2, la règle d'évaluation
$ \int_a^b f = F(b) - F(a) $ pour toute primitive $ F $, en utilisant le théorème
de constance de CAL-07.
\item[(c)] Voyez maintenant pourquoi les hypothèses gagnent leur salaire : posez
$ F(x) = x^2 \sin(1/x^2) $ pour $ x\ne 0 $ et $ F(0)=0 $. Montrez que $ F $ est
dérivable partout, et pourtant que $ F' $ n'est pas bornée sur $ (0,1] $, de sorte que
$ \int_0^1 F' $ n'existe pas comme intégrale de Riemann --- la partie 2 peut être en défaut pour une dérivée
qui n'est pas intégrable.
\end{enumerate}

\end{problem}

\noindent Que le problème de la tangente et le problème de l'aire soient inverses l'un de l'autre
est une \textit{découverte}, entrevue géométriquement par Barrow et exploitée systématiquement par
Newton et Leibniz ; la première démonstration rigoureuse pour des intégrandes continues est celle de Cauchy
(1823). La partie (iii) montre que le théorème est véritablement conditionnel --- et la quête de la
\og{} bonne \fg{} intégrale qui le répare mène à Lebesgue (voir CAL-21 et
Analyse réelle (\url{pure-real-analysis.html})).

\begin{refs}
\item Contexte : Théorème fondamental de l'analyse --- Wikipédia (\url{https://fr.wikipedia.org/wiki/Th%C3%A9or%C3%A8me_fondamental_de_l%27analyse})
\item Lecture : Michael Spivak, \textit{Calculus}, ch. 14
\item Cours : MIT OCW 18.01SC Single Variable Calculus (\url{https://ocw.mit.edu/courses/18-01sc-single-variable-calculus-fall-2010/})
\item Vidéo : 3Blue1Brown, Essence of Calculus (\url{https://www.3blue1brown.com/lessons/essence-of-calculus/})
\end{refs}

\bigskip

\begin{problem}[{La longueur d'arc, avec sa mise en place démontrée --- Lycée}]
\textbf{CAL-10.} \textbf{Problème.} Définissez la longueur d'une courbe comme la borne supérieure des
longueurs des lignes polygonales inscrites.
\begin{enumerate}
\item[(a)] Pour $ f $ continûment dérivable, démontrez
la formule
\[ L = \int_a^b \sqrt{1 + f'(x)^2}\,dx \]
à partir de cette définition --- appliquez le théorème des accroissements finis sur chaque sous-intervalle d'une subdivision et
reconnaissez la somme de Riemann obtenue.
\item[(b)] Servez-vous-en pour calculer la longueur exacte de
$ y = \tfrac{2}{3}x^{3/2} $ de $ x=0 $ à $ x=3 $.
\item[(c)] Défi : calculez la
longueur d'arc de la parabole $ y = x^2/2 $ de $ x=0 $ à $ x=1 $ (l'intégrande
$ \sqrt{1+x^2} $ vous donnera du fil à retordre).
\end{enumerate}

\end{problem}

\noindent Descartes pensait que le rapport entre lignes courbes et lignes droites ne pourrait
jamais être connu exactement. Il avait tort : en 1657 William Neile a rectifié la parabole
semi-cubique $ y^2 = x^3 $ --- essentiellement la courbe de (ii), la première courbe algébrique dont la longueur
fut trouvée exactement --- et Wren a rectifié la cycloïde peu après. L'étape du théorème des accroissements finis en (i) est
le cœur du problème : la formule est habituellement affirmée avec un geste de la main vers des
\og{} triangles rectangles infinitésimaux \fg{}, et vous transformez ici ce geste en démonstration.

\begin{refs}
\item Contexte : Longueur d'un arc --- Wikipédia (\url{https://fr.wikipedia.org/wiki/Longueur_d%27un_arc})
\item Contexte : Parabole semi-cubique --- Wikipédia (\url{https://fr.wikipedia.org/wiki/Parabole_semi-cubique})
\item Lecture : James Stewart, \textit{Calculus}, ch. 8.1
\end{refs}

\bigskip

\begin{problem}[{La série harmonique diverge --- Lycée}]
\textbf{CAL-11.} \textbf{Problème.} Démontrez que
$ 1 + \frac12 + \frac13 + \frac14 + \cdots $ diverge, deux fois :
\begin{enumerate}
\item[(a)] par le
groupement des termes d'Oresme en blocs dépassant chacun $ \frac12 $ ;
\item[(b)] en comparant la somme
partielle $ H_n = \sum_{k=1}^n \frac1k $ avec $ \int_1^n \frac{dx}{x} $. Démontrez ensuite que
$ \sum \frac{1}{n^2} $ \textit{converge} par comparaison avec la série télescopique
$ \sum \frac{1}{n(n-1)} $. Enfin, à partir de vos encadrements de (ii), démontrez que
$ H_n - \ln n $ est décroissante et minorée, donc convergente. Sa limite est un nombre
$ \gamma \approx 0{,}5772 $ --- retenez-le : il revient en CAL-22 avec un problème ouvert
attaché.
\end{enumerate}

\end{problem}

\noindent Nicole Oresme a démontré la divergence vers 1350 ; la démonstration fut oubliée
puis redécouverte par Mengoli et les Bernoulli trois siècles plus tard. La série diverge
si lentement (les $ 10^{43} $ premiers termes n'atteignent pas 100) qu'aucune expérience n'aurait jamais
pu en révéler la vérité --- un pur triomphe de la démonstration sur l'évidence. La valeur de
$ \sum 1/n^2 $ est le problème de Bâle, dont l'histoire vit sur
L'art de la démonstration (\url{proofs.html}).

\begin{refs}
\item Contexte : Série harmonique --- Wikipédia (\url{https://fr.wikipedia.org/wiki/S%C3%A9rie_harmonique})
\item Lecture : Michael Spivak, \textit{Calculus}, chapitres sur les séries infinies
\item Cours : MIT OCW 18.01SC Single Variable Calculus (\url{https://ocw.mit.edu/courses/18-01sc-single-variable-calculus-fall-2010/})
\end{refs}

\bigskip

\begin{problem}[{Les critères de convergence sont des théorèmes --- Lycée}]
\textbf{CAL-12.} \textbf{Problème.} Énoncez et démontrez :
\begin{enumerate}
\item[(a)] le critère de comparaison ;
\item[(b)] la
règle de d'Alembert (par comparaison avec une série géométrique) ;
\item[(c)] le critère des séries alternées,
avec la majoration de l'erreur $ |S - S_n| \le a_{n+1} $ ;
\item[(d)] le test de condensation
de Cauchy : pour $ a_n \ge 0 $ décroissante, $ \sum a_n $ converge si et seulement si
$ \sum 2^k a_{2^k} $ converge. Décidez ensuite, avec démonstration, de la convergence de
\[ \sum_{n\ge 1} \frac{n!}{n^n}, \qquad \sum_{n\ge 1} \frac{(-1)^n}{\sqrt{n}},
\qquad \sum_{n\ge 2} \frac{1}{n \ln n}. \]
\end{enumerate}

\end{problem}

\noindent Les critères ont été assemblés pièce à pièce --- la règle du rapport par d'Alembert,
le critère des séries alternées par Leibniz dans sa correspondance, la condensation par Cauchy, dont le
\textit{Cours d'analyse} (1821) a traité la convergence systématiquement pour la première fois, après un siècle
de manipulations confiantes de séries divergentes qui avaient produit merveilles et absurdités.
Chaque critère est une petite machine ; ici vous ouvrez le boîtier et voyez l'argument de comparaison
à l'intérieur de chacun.

\begin{refs}
\item Contexte : Règle de d'Alembert --- Wikipédia (\url{https://fr.wikipedia.org/wiki/R%C3%A8gle_de_d%27Alembert})
\item Contexte : Critère de convergence des séries alternées --- Wikipédia (\url{https://fr.wikipedia.org/wiki/Crit%C3%A8re_de_convergence_des_s%C3%A9ries_altern%C3%A9es})
\item Contexte : Test de condensation de Cauchy --- Wikipédia (\url{https://fr.wikipedia.org/wiki/Test_de_condensation_de_Cauchy})
\item Lecture : Tom M. Apostol, \textit{Calculus}, vol. 1
\end{refs}

\bigskip

\begin{problem}[{Approcher e, avec un certificat --- Lycée}]
\textbf{CAL-13.} \textbf{Problème.} Énoncez le théorème de Taylor avec le reste sous forme de
Lagrange :
\[ f(x) = \sum_{k=0}^{n} \frac{f^{(k)}(0)}{k!}x^k + \frac{f^{(n+1)}(c)}{(n+1)!}x^{n+1}
\quad\text{pour un certain } c \text{ entre } 0 \text{ et } x. \]
(i) Démontrez d'abord, par comparaison avec une série géométrique, que $ e < 3 $. (ii) À l'aide
du polynôme de Taylor de $ e^x $ en $ x=1 $, trouvez le plus petit $ n $ pour lequel vous
pouvez \textit{démontrer} que l'erreur d'approximation est inférieure à $ 10^{-3} $, et calculez $ e $ avec
cette précision à la main --- l'enjeu n'est pas les décimales mais la majoration certifiée de l'erreur.
(iii) Défi : poussez le même développement plus loin et démontrez que $ e $ est irrationnel
(supposez $ e = p/q $ et multipliez par $ q! $).
\end{problem}

\noindent Brook Taylor a publié sa série en 1715 ; Lagrange a fourni le
reste qui transforme une série formelle en garantie numérique. Euler a calculé $ e $ avec
de nombreuses décimales exactement par la méthode de (ii). La démonstration d'irrationalité de (iii) est celle de Fourier
(1815) et c'est l'un des plus beaux arguments en deux lignes des mathématiques une fois trouvé ---
le trouver est le problème.

\begin{refs}
\item Contexte : Théorème de Taylor --- Wikipédia (\url{https://fr.wikipedia.org/wiki/Th%C3%A9or%C3%A8me_de_Taylor})
\item Contexte : e (constante mathématique) --- Wikipédia (\url{https://fr.wikipedia.org/wiki/E_%28nombre%29})
\item Lecture : Michael Spivak, \textit{Calculus}, ch. 20 (approximation par des fonctions polynomiales)
\item Vidéo : 3Blue1Brown, les séries de Taylor (\url{https://www.3blue1brown.com/lessons/taylor-series-geometric-view/})
\end{refs}

\bigskip

\begin{problem}[{L'Hôpital, justifié --- Lycée}]
\textbf{CAL-14.} \textbf{Problème.}
\begin{enumerate}
\item[(a)] Énoncez et démontrez le théorème des accroissements
finis généralisé de Cauchy : pour $ f, g $ continues sur $ [a,b] $ et dérivables sur $ (a,b) $,
il existe $ c $ tel que $ (f(b)-f(a))\,g'(c) = (g(b)-g(a))\,f'(c) $.
\item[(b)] Servez-vous-en pour démontrer
la règle de L'Hôpital dans le cas $ 0/0 $ lorsque $ x\to a $, avec des
hypothèses soignées.
\item[(c)] Calculez $ \lim_{x\to 0} \frac{1-\cos x}{x^2} $ avec la règle, puis
de nouveau sans elle.
\item[(d)] Montrez que la règle est à sens unique : exhibez $ f, g $ tels que
$ \lim_{x\to\infty} f(x)/g(x) $ existe alors que $ \lim f'(x)/g'(x) $ n'existe pas (essayez
$ f(x) = x + \sin x $, $ g(x) = x $).
\item[(e)] Expliquez pourquoi appliquer la règle à
$ \lim_{x\to 0} \frac{\sin x}{x} $ est circulaire si vous dérivez $ \sin $ de la
manière habituelle.
\end{enumerate}

\end{problem}

\noindent La règle est apparue dans le tout premier manuel d'analyse imprimé ---
l'\textit{Analyse des infiniment petits} de L'Hôpital (1696) --- mais elle avait été découverte par
Jean Bernoulli, que le marquis s'était attaché par contrat pour l'exclusivité de ses
découvertes mathématiques. Bernoulli s'en est plaint amèrement après la mort de L'Hôpital ;
la correspondance a survécu et tranche l'affaire. Le contenu mathématique du
problème est le point (i) : le théorème des accroissements finis de Cauchy est le moteur honnête sous le raccourci populaire.

\begin{refs}
\item Contexte : Règle de L'Hôpital --- Wikipédia (\url{https://fr.wikipedia.org/wiki/R%C3%A8gle_de_L%27H%C3%B4pital})
\item Lecture : Michael Spivak, \textit{Calculus}, ch. 11 et exercices
\item Cours : MIT OCW 18.01SC Single Variable Calculus (\url{https://ocw.mit.edu/courses/18-01sc-single-variable-calculus-fall-2010/})
\end{refs}

\bigskip

\begin{problem}[{L'aire sous la courbe en cloche --- un défi --- Lycée}]
\textbf{CAL-15.} \textbf{Problème.} Démontrez que
\[ \int_{-\infty}^{\infty} e^{-x^2}\,dx = \sqrt{\pi}. \]
Démontrez d'abord que l'intégrale impropre converge (comparez $ e^{-x^2} $ avec
$ e^{-|x|} $ pour $ |x|\ge 1 $). Évaluez-la ensuite. Avertissement loyal, qui relève du contexte et
non de l'indice : $ e^{-x^2} $ n'a pas de primitive élémentaire (c'est un théorème --- voir
CAL-23), de sorte qu'aucune chasse à la primitive ne peut aboutir ; il faut quelque chose de
structurellement différent.
\end{problem}

\noindent Cette intégrale est l'aire totale sous la courbe normale, la raison pour laquelle le
facteur $\surd$$\pi$ hante les probabilités. De Moivre a rencontré la courbe en 1733 en approchant des
comptages de pile ou face ; Laplace a évalué l'intégrale le premier (années 1770), et Gauss a attaché son nom à la
loi en 1809. Lord Kelvin aurait dit à une classe qu'un mathématicien est quelqu'un
pour qui cette identité est aussi évidente que deux fois deux font quatre. Elle n'est pas évidente. Elle est
en revanche trouvable --- et inoubliable une fois trouvée. Sa signification probabiliste est reprise dans
Probabilités (\url{applied-probability.html}) (PRB-14).

\begin{refs}
\item Contexte : Intégrale de Gauss --- Wikipédia (\url{https://fr.wikipedia.org/wiki/Int%C3%A9grale_de_Gauss})
\item Contexte : Loi normale --- Wikipédia (\url{https://fr.wikipedia.org/wiki/Loi_normale})
\item Lecture : Michael Spivak, \textit{Calculus} (intégrales impropres)
\end{refs}

\bigskip

\begin{problem}[{Une limite sans rien à simplifier --- Lycée, short}]
\textbf{CAL-25.} \textbf{Problème.} Démontrez, directement à partir de la définition de la dérivée, que la dérivée de $ f(x) = \sqrt{x} $ en un point $ a > 0 $ vaut $ 1/(2\sqrt{a}) $. Précisez exactement où vous utilisez que $ a \neq 0 $.
\end{problem}

\noindent Le taux d'accroissement d'une racine carrée ne se simplifie pas par factorisation ; il faut multiplier par la quantité conjuguée, une astuce qui revient constamment en analyse. La seconde phrase est le véritable contenu --- une formule de dérivée qui échoue silencieusement en une extrémité est le genre de chose qui devient plus tard une erreur sérieuse.

\begin{refs}
\item Contexte : Taux d'accroissement --- Wikipédia (en anglais) (\url{https://en.wikipedia.org/wiki/Difference_quotient})
\end{refs}

\bigskip

\begin{problem}[{Une série qui se télescope --- Lycée, short}]
\textbf{CAL-26.} \textbf{Problème.} Démontrez que
\[ \sum_{n=1}^{\infty} \frac{1}{n(n+1)} = 1 \]
en trouvant une formule pour la somme partielle $ \sum_{n=1}^{N} \frac{1}{n(n+1)} $ puis en passant à la limite. Justifiez le passage à la limite.
\end{problem}

\noindent Le télescopage est la seule technique de sommation qui donne une réponse exacte sans aucune machinerie, et voici son cas le plus net. Comparez-le à la série harmonique, qui lui ressemble et diverge : remarquer pourquoi l'une s'effondre et l'autre non, c'est commencer à comprendre la convergence.

\begin{refs}
\item Contexte : Somme télescopique --- Wikipédia (\url{https://fr.wikipedia.org/wiki/Somme_t%C3%A9lescopique})
\end{refs}

\bigskip

\begin{problem}[{La boîte de conserve la plus économe --- Lycée, short}]
\textbf{CAL-27.} \textbf{Problème.} Une boîte cylindrique doit contenir un volume fixé $ V $. Déterminez, en démontrant que votre réponse est bien un minimum et non un simple point critique, le rayon et la hauteur qui minimisent l'aire totale de sa surface.
\end{problem}

\noindent C'est le problème d'optimisation classique, et presque tous ceux qui le résolvent s'arrêtent après avoir annulé la dérivée. L'obligation de démonstration est justement le cœur du sujet : une dérivée qui s'annule repère des candidats, et il faut quelque chose de plus --- la dérivée seconde, une étude de signe, ou un argument de compacité --- avant que le mot \og{} minimum \fg{} soit honnête.

\begin{refs}
\item Contexte : Optimisation (mathématiques) --- Wikipédia (\url{https://fr.wikipedia.org/wiki/Optimisation_%28math%C3%A9matiques%29})
\item Lecture : Michael Spivak, \textit{Calculus}, ch. 11
\end{refs}

\bigskip

\begin{problem}[{Le changement de variable, démontré --- Lycée}]
\textbf{CAL-34.} \textbf{Problème.} La règle de substitution dit que si $ g $ est continûment dérivable sur
$ [a,b] $ et si $ f $ est continue sur un intervalle contenant les valeurs de $ g $, alors
\[ \int_a^b f\big(g(x)\big)\,g'(x)\,dx \;=\; \int_{g(a)}^{g(b)} f(u)\,du. \]
\begin{enumerate}
\item[(a)] Démontrez-la à partir de la dérivation des fonctions composées et du théorème fondamental de l'analyse, et dites exactement où
chaque hypothèse sert. Notez que $ g $ n'est \textit{pas} supposée injective, et expliquez
pourquoi le théorème y survit.
\item[(b)] Servez-vous-en pour évaluer $ \int_0^{\pi} \sin^3 x \,dx $ et $ \int_0^1 x\sqrt{1-x^2}\,dx $,
et vérifiez la seconde réponse indépendamment en interprétant l'intégrale comme une aire.
\item[(c)] Voici maintenant le sens qui, lui, exige l'injectivité : énoncez et démontrez la formule de changement
de variable sous la forme $ \int_c^d f(u)\,du = \int_{g^{-1}(c)}^{g^{-1}(d)} f(g(x))g'(x)\,dx $,
et exhibez un $ g $ non monotone pour lequel cette lecture inverse est en défaut.
\item[(d)] Deux pièges. Premièrement, dans $ \int_{-1}^{1} \frac{dx}{1+x^2} $ substituez $ x = 1/u $ ; vous
obtiendrez l'équation $ I = -I $, donc $ I = 0 $, pour l'intégrale d'une fonction strictement
positive. Deuxièmement, appliquez la substitution en tangente de l'arc moitié $ t = \tan(x/2) $ à
$ \int_0^{2\pi} \frac{dx}{2+\cos x} $ et regardez-la renvoyer $ 0 $ elle aussi. Dans chaque cas, trouvez
l'hypothèse de (a) qui a été violée, et réparez le calcul.
\end{enumerate}

\end{problem}

\noindent La notation $ dx $ de Leibniz a été conçue pour que le changement de variable ait l'air
d'une simplification de symboles, et c'est la principale raison pour laquelle sa notation l'a emporté sur celle de Newton. Mais la
simplification est un théorème, et la partie (d) montre ce qui arrive quand on lui fait confiance comme à une comptabilité :
les deux sophismes sont de vieux favoris des salles de classe, et tous deux viennent d'une substitution qui n'est pas
dérivable --- voire pas même définie --- quelque part à l'intérieur de l'intervalle. La substitution en arc moitié
de (d) est habituellement appelée substitution de Weierstrass, une attribution que personne n'a
jamais su documenter ; elle apparaît chez Euler.

\begin{refs}
\item Contexte : Intégration par changement de variable --- Wikipédia (\url{https://fr.wikipedia.org/wiki/Int%C3%A9gration_par_changement_de_variable})
\item Contexte : Substitution en tangente de l'arc moitié --- Wikipédia (en anglais) (\url{https://en.wikipedia.org/wiki/Tangent_half-angle_substitution})
\item Lecture : Tom M. Apostol, \textit{Calculus}, vol. 1 (le théorème de substitution)
\item Cours : MIT OCW 18.01SC Single Variable Calculus (\url{https://ocw.mit.edu/courses/18-01sc-single-variable-calculus-fall-2010/})
\end{refs}

\bigskip

\begin{problem}[{Aires en polaires et arche de la cycloïde --- Lycée}]
\textbf{CAL-35.} \textbf{Problème.} Une courbe peut être donnée sous forme polaire $ r = r(\theta) $, ou paramétriquement par
$ \big(x(t), y(t)\big) $.
\begin{enumerate}
\item[(a)] Démontrez la formule de l'aire en polaires
\[ A \;=\; \frac{1}{2}\int_{\alpha}^{\beta} r(\theta)^2 \, d\theta \]
pour $ r \ge 0 $ continue : approchez par des secteurs circulaires plutôt que par des rectangles, écrivez
les sommes supérieures et inférieures, et montrez qu'elles convergent vers la même limite.
\item[(b)] Calculez, avec démonstration, l'aire enfermée par la cardioïde $ r = 1+\cos\theta $ et l'aire
d'un pétale de la rosace $ r = \cos 3\theta $ ; expliquez ensuite pourquoi cette rosace a trois pétales
tandis que $ r = \cos 2\theta $ en a quatre.
\item[(c)] Calculez l'aire balayée par la spirale d'Archimède $ r = \theta $ pour
$ 0 \le \theta \le 2\pi $, et comparez-la à l'aire du cercle $ r = 2\pi $ --- la
comparaison qu'Archimède a démontrée dans \textit{Des spirales} sans aucune intégrale.
\item[(d)] Démontrez la formule de longueur d'arc en paramétriques
$ L = \int_{t_0}^{t_1} \sqrt{x'(t)^2 + y'(t)^2}\,dt $ pour $ x, y $ continûment
dérivables, et servez-vous-en, avec la formule de l'aire
$ A = \int y \, dx $, pour trouver la longueur d'une arche de la cycloïde
$ x = r(t-\sin t) $, $ y = r(1-\cos t) $, $ 0 \le t \le 2\pi $, et l'aire entre cette
arche et le sol.
\item[(e)] L'arche possède un point de rebroussement : montrez que $ x'(0) = y'(0) = 0 $, de sorte que la courbe n'y a pas
de direction tangente, et expliquez pourquoi la formule de longueur d'arc de (d) reste néanmoins valable.
\end{enumerate}

\end{problem}

\noindent Jacob Bernoulli a publié le premier usage systématique des coordonnées polaires dans les
\textit{Acta Eruditorum} en 1691 ; les rosaces ont été nommées et étudiées par Guido Grandi dans les
années 1720. La cycloïde fut l'obsession du dix-septième siècle --- \og{} l'Hélène des
géomètres \fg{}, ainsi nommée pour les querelles qu'elle a provoquées. Roberval a trouvé l'aire sous une arche en
1634, Christopher Wren a rectifié l'arche exactement en 1658, Huygens a découvert en 1659 qu'elle est
la tautochrone et a bâti des horloges à pendule sur ce fait, et le défi lancé par Jean Bernoulli en 1696
l'a révélée brachistochrone, la courbe de descente la plus rapide. Tout cela a été fait
avec la machinerie que vous démontrez ici, ou avec moins.

\begin{refs}
\item Contexte : Coordonnées polaires --- Wikipédia (\url{https://fr.wikipedia.org/wiki/Coordonn%C3%A9es_polaires})
\item Contexte : Cycloïde --- Wikipédia (\url{https://fr.wikipedia.org/wiki/Cyclo%C3%AFde})
\item Contexte : Courbe tautochrone --- Wikipédia (\url{https://fr.wikipedia.org/wiki/Courbe_tautochrone})
\item Lecture : James Stewart, \textit{Calculus}, ch. 10 (équations paramétriques et coordonnées polaires)
\item Lecture : E. Hairer \& G. Wanner, \textit{Analysis by Its History} (Springer)
\end{refs}

\bigskip

\begin{problem}[{Les fonctions qui mesurent une chaîne suspendue --- Lycée, short}]
\textbf{CAL-36.} \textbf{Problème.} Posez
$ \cosh x = \frac{e^x + e^{-x}}{2} $ et $ \sinh x = \frac{e^x - e^{-x}}{2} $. Démontrez que
$ \cosh^2 x - \sinh^2 x = 1 $, que $ \frac{d}{dx}\cosh x = \sinh x $, et que la longueur
d'arc de la courbe $ y = \cosh x $ de $ x = 0 $ à $ x = a $ vaut exactement
$ \sinh a $.
\end{problem}

\noindent Vincenzo Riccati a introduit ces combinaisons dans les années 1760 et Lambert
les a fait entrer dans l'usage courant ; l'identité de la première tâche est la raison pour laquelle on les dit hyperboliques
--- le point $ (\cosh t, \sinh t) $ parcourt l'hyperbole $ x^2 - y^2 = 1 $ exactement comme
$ (\cos t, \sin t) $ parcourt le cercle. La dernière tâche est le petit miracle derrière CAL-37 :
l'unique courbe dont l'intégrale de longueur d'arc s'effondre est celle selon laquelle une chaîne pend réellement.

\begin{refs}
\item Contexte : Fonction hyperbolique --- Wikipédia (\url{https://fr.wikipedia.org/wiki/Fonction_hyperbolique})
\item Contexte : Chaînette --- Wikipédia (\url{https://fr.wikipedia.org/wiki/Cha%C3%AEnette})
\end{refs}

\bigskip


\section*{Licence}

\begin{problem}[{Les multiplicateurs de Lagrange, avec justification --- Licence}]
\textbf{CAL-16.} \textbf{Problème.}
\begin{enumerate}
\item[(a)] Utilisez les multiplicateurs de Lagrange pour démontrer que parmi tous les
pavés rectangles d'aire totale fixée, le cube a le plus grand volume.
\item[(b)] Démontrez le théorème que vous avez utilisé : si $ f $ et $ g $ sont $ C^1 $, si $ a $ est un
extremum local de $ f $ restreinte à la ligne de niveau $ \{g = c\} $, et si
$ \nabla g(a) \ne 0 $, alors $ \nabla f(a) = \lambda \nabla g(a) $ pour un certain
$ \lambda\in\mathbb{R} $. (Dérivez $ f $ le long de courbes tracées dans la ligne de niveau,
et expliquez --- via le théorème des fonctions implicites --- pourquoi il existe assez de telles courbes.)
\item[(c)] Montrez que l'hypothèse $ \nabla g \ne 0 $ n'est pas décorative : minimisez
$ f(x,y) = x $ sur la courbe $ x^3 = y^2 $ et vérifiez qu'aucun $ \lambda $ ne convient au
point minimisant.
\end{enumerate}

\end{problem}

\noindent Lagrange a introduit les multiplicateurs dans la \textit{Mécanique analytique}
(1788) pour traiter les équilibres sous contrainte sans éliminer de variables. La méthode est
enseignée partout et justifiée presque nulle part ; le point de rebroussement de (iii), où l'ensemble de
contrainte n'a pas de direction tangente, est la mise en garde standard, et le multiplicateur lui-même
prend une vie propre en économie (les prix implicites) et en physique.

\begin{refs}
\item Contexte : Multiplicateur de Lagrange --- Wikipédia (\url{https://fr.wikipedia.org/wiki/Multiplicateur_de_Lagrange})
\item Lecture : Tom M. Apostol, \textit{Calculus}, vol. 2
\item Cours : MIT OCW 18.02SC Multivariable Calculus (\url{https://ocw.mit.edu/courses/18-02sc-multivariable-calculus-fall-2010/})
\end{refs}

\bigskip

\begin{problem}[{Le théorème de Green à l'œuvre --- Licence}]
\textbf{CAL-17.} \textbf{Problème.} Énoncez le théorème de Green pour une courbe fermée simple
$ C $, orientée positivement et de classe $ C^1 $ par morceaux, bordant un domaine $ D $ :
\[ \oint_C P\,dx + Q\,dy \;=\; \iint_D \Big( \frac{\partial Q}{\partial x} -
\frac{\partial P}{\partial y} \Big)\,dA. \]
(i) Démontrez-le pour les domaines qui sont simultanément de type I et de type II (graphes dans les deux
directions), en ramenant chaque intégrale double au théorème fondamental de l'analyse.
(ii) Déduisez-en la formule de l'aire $ \mathrm{Aire}(D) = \tfrac12 \oint_C (x\,dy - y\,dx) $.
(iii) À partir de (ii), établissez la formule du lacet pour l'aire d'un polygone de sommets
$ (x_1,y_1),\dots,(x_n,y_n) $. (iv) Expliquez, la formule en main, comment un
planimètre --- un bras mécanique articulé --- mesure l'aire d'un domaine quelconque en se contentant d'en
parcourir le bord.
\end{problem}

\noindent George Green était le fils autodidacte d'un meunier de Nottingham ; il a
publié par souscription, en 1828, l'essai contenant le cercle d'idées du théorème,
devant un public d'une cinquantaine de personnes ; Kelvin a redécouvert et défendu ce travail en 1846, quatre
ans après la mort de Green. Le théorème est le cas plan du théorème de Stokes, la
grande unification des théorèmes fondamentaux en dimension un et en dimension supérieure.

\begin{refs}
\item Contexte : Théorème de Green --- Wikipédia (\url{https://fr.wikipedia.org/wiki/Th%C3%A9or%C3%A8me_de_Green})
\item Contexte : Planimètre --- Wikipédia (\url{https://fr.wikipedia.org/wiki/Planim%C3%A8tre})
\item Lecture : Tom M. Apostol, \textit{Calculus}, vol. 2
\item Cours : MIT OCW 18.02SC Multivariable Calculus (\url{https://ocw.mit.edu/courses/18-02sc-multivariable-calculus-fall-2010/})
\end{refs}

\bigskip

\begin{problem}[{Quand les dérivées partielles croisées commutent --- Licence}]
\textbf{CAL-18.} \textbf{Problème.} (i) Démontrez le théorème de Schwarz : si les dérivées partielles
secondes croisées $ f_{xy} $ et $ f_{yx} $ existent et sont continues au voisinage d'un point, elles y sont
égales. (Appliquez deux fois le théorème des accroissements finis à la différence seconde
$ f(x+h,y+k) - f(x+h,y) - f(x,y+k) + f(x,y) $.) (ii) Calculez ensuite les deux dérivées partielles croisées
à l'origine pour
\[ f(x,y) = \frac{xy(x^2-y^2)}{x^2+y^2}, \qquad f(0,0)=0, \]
et montrez qu'elles existent mais ne sont \textit{pas} égales --- localisez exactement quelle hypothèse de (i)
tombe en défaut.
\end{problem}

\noindent Euler et Clairaut utilisaient librement cette symétrie dans les années 1730-40 ; la démonstration
publiée par Clairaut était erronée, et la première démonstration correcte sous des hypothèses propres est
celle de H. A. Schwarz (1873). Le contre-exemple de (ii) est celui de Peano (1884). La morale
traverse toute l'analyse : intervertir deux passages à la limite est un théorème à chaque
fois, jamais un réflexe --- un thème poursuivi sur
Analyse réelle (\url{pure-real-analysis.html}).

\begin{refs}
\item Contexte : Symétrie des dérivées secondes --- Wikipédia (\url{https://fr.wikipedia.org/wiki/Th%C3%A9or%C3%A8me_de_Schwarz})
\item Lecture : Tom M. Apostol, \textit{Calculus}, vol. 2
\item Cours : MIT OCW 18.02SC Multivariable Calculus (\url{https://ocw.mit.edu/courses/18-02sc-multivariable-calculus-fall-2010/})
\end{refs}

\bigskip

\begin{problem}[{Dériver sous le signe intégral --- Licence}]
\textbf{CAL-19.} \textbf{Problème.} (i) Démontrez la règle de Leibniz : si $ f(x,t) $ et
$ \partial f/\partial t $ sont continues sur $ [a,b]\times[c,d] $, alors
$ F(t) = \int_a^b f(x,t)\,dx $ est dérivable et
$ F'(t) = \int_a^b \frac{\partial f}{\partial t}(x,t)\,dx $. (Utilisez la continuité uniforme de
$ \partial f/\partial t $ sur le rectangle compact.) (ii) Servez-vous de la technique pour évaluer
\[ \int_0^1 \frac{x^b - 1}{\ln x}\,dx \;(b>0) \qquad\text{et}\qquad
\int_0^\infty e^{-x}\,\frac{\sin x}{x}\,dx ; \]
dans le cas impropre, énoncez et justifiez l'hypothèse supplémentaire dont vous avez besoin. (iii) Trouvez, ou
construisez, un exemple où l'interversion aveugle de $ d/dt $ et de $ \int $ donne une réponse
fausse, et identifiez l'hypothèse en défaut.
\end{problem}

\noindent Feynman a appris l'astuce, comme chacun sait, dans l'\textit{Advanced
Calculus} de Woods au lycée et \og{} s'est servi encore et encore de ce fichu outil \fg{},
résolvant des intégrales qui bloquaient des collègues formés aux seules méthodes de contour. C'est un parfait
spécimen du thème de l'interversion des limites : fabuleusement efficace, et valide seulement parce
qu'un théorème le soutient.

\begin{refs}
\item Contexte : Dérivation sous le signe intégral --- Wikipédia (en anglais) (\url{https://en.wikipedia.org/wiki/Leibniz_integral_rule})
\item Lecture : Richard P. Feynman, \textit{Surely You're Joking, Mr. Feynman!}, \og{} A Different Box of Tools \fg{}
\item Lecture : Tom M. Apostol, \textit{Calculus}, vol. 2
\end{refs}

\bigskip

\begin{problem}[{Des dérivées partielles croisées qui refusent de commuter --- Licence, short}]
\textbf{CAL-28.} \textbf{Problème.} Pour
\[ f(x,y) = \frac{xy(x^2 - y^2)}{x^2 + y^2} \quad (x,y) \neq (0,0), \qquad f(0,0) = 0, \]
montrez que $ f_{xy}(0,0) \neq f_{yx}(0,0) $, et identifiez précisément quelle hypothèse du théorème de Clairaut tombe en défaut.
\end{problem}

\noindent Tout cours de calcul différentiel à plusieurs variables énonce que les dérivées partielles croisées commutent, et la plupart des étudiants ne rencontrent jamais le contre-exemple, ce qui donne au théorème des allures de formalité plutôt que de résultat au contenu réel. Localiser l'hypothèse en défaut --- la continuité des dérivées partielles secondes au voisinage de l'origine --- est ce qui convertit une règle apprise par cœur en connaissance.

\begin{refs}
\item Contexte : Symétrie des dérivées secondes --- Wikipédia (\url{https://fr.wikipedia.org/wiki/Th%C3%A9or%C3%A8me_de_Schwarz})
\end{refs}

\bigskip

\begin{problem}[{Un critère intégral qu'il faut justifier --- Licence, short}]
\textbf{CAL-29.} \textbf{Problème.} Déterminez pour quels réels $ p $ la série $ \sum_{n=2}^{\infty} \frac{1}{n (\log n)^{p}} $ converge, et démontrez votre réponse à l'aide de la comparaison série-intégrale --- en énonçant les hypothèses de ce critère et en vérifiant que chacune est satisfaite ici.
\end{problem}

\noindent Cette famille se situe exactement à la frontière où les critères de comparaison plus simples ne donnent aucune information, ce qui en fait le terrain d'essai standard de la comparaison série-intégrale. La consigne de vérifier les hypothèses est délibérée : la monotonie tombe en défaut pour les petits $ n $, et remarquer que le critère s'applique tout de même à partir d'un certain rang fait partie de l'exercice.

\begin{refs}
\item Contexte : Comparaison série-intégrale --- Wikipédia (\url{https://fr.wikipedia.org/wiki/Comparaison_s%C3%A9rie-int%C3%A9grale})
\end{refs}

\bigskip

\begin{problem}[{Pourquoi une chaîne suspendue n'est pas une parabole --- Licence}]
\textbf{CAL-37.} \textbf{Problème.} Une chaîne parfaitement souple et inextensible, de masse linéique constante
$ \rho $, pend en équilibre sous la gravité $ g $, soutenue à ses deux extrémités. Soit $ y(x) $
sa forme et soit $ T_0 $ la composante horizontale de la tension, qui est la même en
tout point.
\begin{enumerate}
\item[(a)] Équilibrez les forces sur le morceau de chaîne compris entre le point le plus bas et un point général, et
établissez
\[ y'' \;=\; \frac{\rho g}{T_0}\,\sqrt{1 + (y')^2}. \]
Dites en une phrase ce que la racine carrée fait là, et en quoi l'équation différerait si
le poids était réparti uniformément le long de l'\textit{horizontale} au lieu de l'être le long de la chaîne.
\item[(b)] Résolvez l'équation : substituez $ p = y' $, séparez les variables, et démontrez que toute
solution est de la forme $ y = a\cosh\frac{x-c}{a} + d $ avec $ a = T_0/(\rho g) $. Démontrez ensuite
que les constantes sont déterminées par les deux points d'appui et la longueur de la chaîne.
\item[(c)] Démontrez que le câble d'un pont suspendu, qui porte un tablier de poids constant par
mètre horizontal et dont le poids propre est négligeable, pend au contraire selon une parabole. Comparez ensuite
quantitativement les deux courbes pour une portée dont la flèche vaut un dixième de la largeur, et décidez si
la différence est visible.
\item[(d)] Démontrez que la chaîne minimise l'énergie potentielle : parmi toutes les courbes de longueur fixée joignant
deux points fixés, celle dont le centre de masse est le plus bas satisfait l'équation de (a). Posez le
problème comme la minimisation de $ \int y\sqrt{1+(y')^2}\,dx $ sous la contrainte
$ \int \sqrt{1+(y')^2}\,dx = L $, introduisez un multiplicateur de Lagrange, et établissez
l'équation d'Euler--Lagrange.
\end{enumerate}

\end{problem}

\noindent Galilée a affirmé dans les \textit{Discorsi} (1638) qu'une chaîne suspendue prend la
forme d'une parabole ; Joachim Jungius a montré qu'il n'en est rien, dans un travail publié en 1669. Jakob
Bernoulli a lancé le problème comme défi public dans les \textit{Acta Eruditorum} en 1690, et en 1691
trois solutions sont arrivées --- de son frère Johann, de Leibniz, et de Huygens, qui a forgé le
nom de \textit{catenaria}. Johann n'a jamais cessé de savourer que Jakob ait échoué à résoudre son propre
problème. Le bénéfice pour l'ingénierie est venu de Robert Hooke, qui a publié en 1675 sous forme d'anagramme le
principe selon lequel une arche doit être une chaîne renversée ; Gaudí suspendait des ficelles lestées à l'envers pour
concevoir la Sagrada Família, et la Gateway Arch de Saint-Louis est une chaînette aplatie.

\begin{refs}
\item Contexte : Chaînette --- Wikipédia (\url{https://fr.wikipedia.org/wiki/Cha%C3%AEnette})
\item Contexte : Calcul des variations --- Wikipédia (\url{https://fr.wikipedia.org/wiki/Calcul_des_variations})
\item Lecture : E. Hairer \& G. Wanner, \textit{Analysis by Its History} (Springer)
\item Voir aussi : Équations différentielles (\url{applied-ode.html}), Physique mathématique (\url{applied-physics.html})
\end{refs}

\bigskip

\begin{problem}[{Le produit de Wallis, et la constante de la formule de Stirling --- Licence}]
\textbf{CAL-38.} \textbf{Problème.} Soit $ I_n = \int_0^{\pi/2} \sin^n x \, dx $ pour les entiers $ n \ge 0 $.
\begin{enumerate}
\item[(a)] Démontrez par intégration par parties la formule de réduction $ I_n = \frac{n-1}{n} I_{n-2} $ pour
$ n \ge 2 $, et écrivez des expressions closes de $ I_{2m} $ et $ I_{2m+1} $ en termes
de factorielles.
\item[(b)] Démontrez que $ (I_n) $ est strictement décroissante et que $ I_{n}/I_{n-2} \to 1 $. Déduisez
de l'encadrement $ I_{2m+1} \le I_{2m} \le I_{2m-1} $ le produit de Wallis
\[ \frac{\pi}{2} \;=\; \prod_{n=1}^{\infty} \frac{4n^2}{4n^2-1}
\;=\; \frac{2}{1}\cdot\frac{2}{3}\cdot\frac{4}{3}\cdot\frac{4}{5}\cdot\frac{6}{5}\cdot\frac{6}{7}\cdots \]
en énonçant soigneusement ce que signifie la convergence d'un produit infini.
\item[(c)] Démontrez qu'il existe une constante $ C > 0 $ telle que
$ n! = C\, n^{n+1/2} e^{-n} \big(1 + o(1)\big) $. (Comparez $ \ln n! = \sum_{k\le n} \ln k $
avec $ \int_1^n \ln x \, dx $, et montrez que les termes d'erreur forment une série convergente.) Servez-vous ensuite
de (b) pour identifier $ C $ exactement.
\item[(d)] Appliquez le résultat au coefficient binomial central $ \binom{2n}{n} $, et démontrez que la
probabilité d'obtenir exactement $ n $ faces en $ 2n $ lancers d'une pièce équilibrée est équivalente à
$ 1/\sqrt{\pi n} $.
\end{enumerate}

\end{problem}

\noindent John Wallis a obtenu son produit dans l'\textit{Arithmetica Infinitorum} (1656) par une
audacieuse interpolation entre exposants entiers --- un raisonnement que Newton a lu étudiant et transformé
en série du binôme. La partie (c) respecte l'ordre historique des événements : Abraham de Moivre
a trouvé la forme de l'asymptotique de la factorielle en étudiant le jeu de pile ou face dans les années 1730 mais n'a
pas su fixer la constante multiplicative, et c'est James Stirling qui l'a identifiée comme
$ \sqrt{2\pi} $, par un argument exactement de ce type. La partie (d) est la graine du théorème central
limite : $\surd$$\pi$ apparaît dans une question sur des pièces parce qu'il est apparu dans une question sur
$ \sin^n $.

\begin{refs}
\item Contexte : Produit de Wallis --- Wikipédia (\url{https://fr.wikipedia.org/wiki/Produit_de_Wallis})
\item Contexte : Formule de Stirling --- Wikipédia (\url{https://fr.wikipedia.org/wiki/Formule_de_Stirling})
\item Lecture : Tom M. Apostol, \textit{Calculus}, vol. 1
\item Voir aussi : Probabilités (\url{applied-probability.html})
\end{refs}

\bigskip

\begin{problem}[{Quelles intégrales impropres convergent --- Licence, short}]
\textbf{CAL-39.} \textbf{Problème.} Déterminez, avec démonstration, exactement quels couples de réels
$ (a,b) $ rendent
\[ \int_0^{\infty} \frac{x^{a}}{1+x^{b}}\,dx \]
convergente. Traitez séparément le comportement en $ 0 $ et en $ \infty $, et justifiez chaque
comparaison que vous employez en énonçant le critère de comparaison pour les intégrales impropres et en vérifiant ses
hypothèses.
\end{problem}

\noindent Presque toute question de convergence sur une intégrale impropre se ramène, après
une estimation, à la famille $ \int x^{p} $ au voisinage d'un mauvais point --- de sorte que ce seul problème à deux
paramètres contient l'essentiel des exemples d'un premier cours. L'intégrale est aussi une
fonction Bêta déguisée, et c'est ainsi qu'Euler l'a rencontrée dans les années 1730 ; découper le domaine en
$ x = 1 $ et comparer sur chaque moitié constitue toute la méthode.

\begin{refs}
\item Contexte : Intégrale impropre --- Wikipédia (\url{https://fr.wikipedia.org/wiki/Int%C3%A9grale_impropre})
\item Contexte : Règles de comparaison --- Wikipédia (\url{https://fr.wikipedia.org/wiki/S%C3%A9rie_convergente})
\end{refs}

\bigskip


\section*{Master}

\begin{problem}[{Un monstre : continue partout, dérivable nulle part --- Master}]
\textbf{CAL-20.} \textbf{Problème.} Construisez une fonction $ f:\mathbb{R}\to\mathbb{R} $
continue en tout point et dérivable en aucun, et démontrez les deux propriétés.
Deux voies standard, l'une ou l'autre acceptable : celle de Weierstrass,
$ f(x) = \sum_{n=0}^{\infty} a^n \cos(b^n \pi x) $ avec $ 0<a<1 $ et $ ab $
assez grand ; ou la construction de van der Waerden--Takagi
$ \sum_{n=0}^{\infty} 2^{-n}\,\mathrm{dist}(2^n x, \mathbb{Z}) $, dont l'argument de
non-dérivabilité est plus élémentaire. Dans les deux cas, énoncez et
démontrez d'abord le critère M de Weierstrass, qui donne la continuité de la somme. Bonus : cherchez et
comprenez (énoncé seulement) le résultat de Banach de 1931 selon lequel, au sens de la catégorie
de Baire, \textit{presque toutes} les fonctions continues sur $ [0,1] $ sont dérivables nulle part.
\end{problem}

\noindent Weierstrass a présenté son exemple à l'Académie de Berlin en 1872,
démolissant le consensus de l'époque selon lequel une fonction continue doit être dérivable
sauf en des points isolés (Bolzano avait une construction analogue vers 1830, non publiée).
Hermite a écrit qu'il se détournait \og{} avec effroi et horreur de cette plaie lamentable des
fonctions sans dérivée \fg{}. Le monstre marque la porte qui mène de l'analyse élémentaire à l'analyse
réelle ; toute la théorie qui le sous-tend vit sur
Analyse réelle (\url{pure-real-analysis.html}).

\begin{refs}
\item Contexte : Fonction de Weierstrass --- Wikipédia (\url{https://fr.wikipedia.org/wiki/Fonction_de_Weierstrass})
\item Lecture : Stephen Abbott, \textit{Understanding Analysis}, ch. 5
\item Cours : MIT OCW 18.100A Real Analysis (\url{https://ocw.mit.edu/courses/18-100a-real-analysis-fall-2020/})
\end{refs}

\bigskip

\begin{problem}[{Quelles fonctions sont intégrables au sens de Riemann ? --- Master}]
\textbf{CAL-21.} \textbf{Problème.}
\begin{enumerate}
\item[(a)] La fonction de Thomae vaut $ 1/q $ en chaque
rationnel $ p/q $ écrit sous forme irréductible et $ 0 $ en chaque irrationnel. Démontrez qu'elle est continue
exactement aux irrationnels, et intégrable au sens de Riemann sur $ [0,1] $ d'intégrale $ 0 $.
\item[(b)] Démontrez que la fonction de Dirichlet ($ 1 $ sur $ \mathbb{Q} $, $ 0 $ ailleurs)
n'est continue nulle part et n'est pas intégrable au sens de Riemann.
\item[(c)] Définissez ce que signifie, pour un ensemble de
réels, être de mesure nulle. Énoncez le critère de Lebesgue --- une fonction bornée sur
$ [a,b] $ est intégrable au sens de Riemann si et seulement si son ensemble de points de discontinuité est de mesure
nulle --- confrontez-le à vos verdicts en (i) et (ii), et démontrez la moitié facile :
les ensembles dénombrables sont de mesure nulle.
\item[(d)] Bonus : montrez qu'aucune fonction
$ \mathbb{R}\to\mathbb{R} $ ne peut être continue exactement aux rationnels.
\end{enumerate}

\end{problem}

\noindent Riemann a défini son intégrale dans sa leçon d'habilitation de 1854 en partie pour
poser précisément cette question : à quel point une fonction intégrable peut-elle être discontinue ? L'exemple de
Thomae (1875) et celui de Dirichlet balisent le territoire, et la thèse de Lebesgue de 1902
a répondu complètement --- en chemin vers une nouvelle intégrale qui répare le défaut trouvé en
CAL-09(iii). Ce problème est le passeport de l'étudiant d'analyse élémentaire vers
Analyse réelle (\url{pure-real-analysis.html}).

\begin{refs}
\item Contexte : Fonction de Thomae --- Wikipédia (\url{https://fr.wikipedia.org/wiki/Fonction_de_Thomae})
\item Contexte : Intégrale de Riemann --- Wikipédia (\url{https://fr.wikipedia.org/wiki/Int%C3%A9grale_de_Riemann})
\item Lecture : Stephen Abbott, \textit{Understanding Analysis}
\item Cours : MIT OCW 18.100A Real Analysis (\url{https://ocw.mit.edu/courses/18-100a-real-analysis-fall-2020/})
\end{refs}

\bigskip

\begin{problem}[{Un théorème qui contient tous les autres --- Master, short}]
\textbf{CAL-30.} \textbf{Problème.} Énoncez le théorème de Stokes généralisé pour une
$ k $-forme lisse $ \omega $ sur une variété orientée à bord,
\[ \int_{M} d\omega = \int_{\partial M} \omega , \]
puis montrez que le théorème fondamental de l'analyse, le théorème de Green et le théorème
de la divergence en sont chacun le cas particulier pour un $ M $ et un $ \omega $ donnés.
\end{problem}

\noindent Quatre théorèmes enseignés en quatre semaines différentes, avec quatre
démonstrations différentes et quatre notations différentes, se révèlent être un seul énoncé sur la façon dont
dérivation et bords échangent leurs places. Le reconnaître est le moment où l'analyse cesse
d'être une liste de techniques et devient une discipline --- et c'est pourquoi les formes différentielles, qui
ressemblent à une abstraction gratuite quand on les introduit, gagnent immédiatement leur salaire.

\begin{refs}
\item Contexte : Théorème de Stokes généralisé --- Wikipédia (\url{https://fr.wikipedia.org/wiki/Th%C3%A9or%C3%A8me_de_Stokes})
\item Lecture : Michael Spivak, \textit{Calculus on Manifolds} --- le classique court traitement
\item Voir aussi : Topologie (\url{pure-topology.html}), Physique mathématique (\url{applied-physics.html})
\end{refs}

\bigskip

\begin{problem}[{Là où les séries de Taylor ne disent la vérité sur rien --- Master, short}]
\textbf{CAL-31.} \textbf{Problème.} Soit $ f(x) = e^{-1/x^2} $ pour $ x \neq 0 $ et $ f(0) = 0 $. Démontrez que $ f $ est indéfiniment dérivable sur $ \mathbb{R} $ avec $ f^{(n)}(0) = 0 $ pour tout $ n $, et déduisez-en que sa série de Taylor en $ 0 $ converge partout mais ne représente $ f $ nulle part sauf à l'origine.
\end{problem}

\noindent Ce seul exemple sépare le lisse de l'analytique et explique pourquoi l'analyse complexe est une discipline différente de l'analyse réelle : sur $ \mathbb{C} $ aucune fonction de ce genre ne peut exister. C'est aussi la graine des fonctions plateau et des partitions de l'unité, et c'est ainsi que l'exemple gagne sa place en géométrie différentielle au lieu de rester une curiosité.

\begin{refs}
\item Contexte : Fonction régulière non analytique --- Wikipédia (\url{https://fr.wikipedia.org/wiki/Fonction_r%C3%A9guli%C3%A8re_non_analytique})
\item Voir aussi : Analyse complexe (\url{pure-complex-analysis.html})
\end{refs}

\bigskip

\begin{problem}[{La fonction exponentielle, définie de quatre manières --- Master}]
\textbf{CAL-40.} \textbf{Problème.} Considérez quatre définitions candidates de la fonction exponentielle sur
$ \mathbb{R} $ :
\[ E(x) = \sum_{n=0}^{\infty} \frac{x^n}{n!}, \qquad
y' = y \text{ avec } y(0)=1, \qquad
L^{-1} \text{ où } L(x) = \int_1^{x} \frac{dt}{t}, \qquad
P(x) = \lim_{n\to\infty}\Big(1+\frac{x}{n}\Big)^n . \]
\begin{enumerate}
\item[(a)] Démontrez que la série converge pour tout réel $ x $, qu'elle peut être dérivée terme
à terme (énoncez et démontrez le théorème qui l'autorise), et que $ E' = E $. Démontrez la
loi d'addition $ E(x+y) = E(x)E(y) $ à partir du produit de Cauchy de deux séries absolument
convergentes.
\item[(b)] Démontrez que le problème de Cauchy a exactement une solution --- considérez
$ \frac{d}{dx}\big( y(x) E(-x) \big) $ --- de sorte que les définitions un et deux coïncident.
\item[(c)] Démontrez que $ L $ est strictement croissante sur $ (0,\infty) $, satisfait
$ L(uv) = L(u)+L(v) $, et envoie $ (0,\infty) $ sur $ \mathbb{R} $ tout entier ; concluez que
$ L^{-1} $ existe, et démontrez qu'elle est égale à $ E $.
\item[(d)] Démontrez que $ P(x) $ existe pour tout $ x $ et vaut $ E(x) $, la convergence étant
uniforme sur les bornés.
\item[(e)] Comparez maintenant les quatre voies en tant qu'outils. Laquelle donne la loi d'addition le plus à bon compte ? Laquelle
démontre $ 2 < e < 3 $ avec le moins de travail ? Lesquelles s'étendent mot pour mot aux arguments
complexes, et lesquelles s'étendent aux matrices carrées ? Répondez à chaque question par une démonstration ou une
obstruction précise.
\end{enumerate}

\end{problem}

\noindent Euler a traité l'exponentielle et le logarithme par des procédés infinis dans
l'\textit{Introductio in analysin infinitorum} (1748), passant de la limite du produit à la
série sans se soucier de savoir laquelle était la définition ; la limite du produit est de lui. L'habitude
moderne de choisir une définition et d'en déduire les autres est une discipline du dix-neuvième siècle, et elle
mérite d'être pratiquée une fois complètement, parce que chaque cours d'analyse ultérieur en choisit silencieusement une différente.
Rudin ouvre \textit{Real and Complex Analysis} par un prologue consacré à cette seule fonction,
qu'il appelle la plus importante des mathématiques.

\begin{refs}
\item Contexte : Caractérisations de la fonction exponentielle --- Wikipédia (en anglais) (\url{https://en.wikipedia.org/wiki/Characterizations_of_the_exponential_function})
\item Contexte : Fonction exponentielle --- Wikipédia (\url{https://fr.wikipedia.org/wiki/Fonction_exponentielle})
\item Lecture : Walter Rudin, \textit{Real and Complex Analysis}, prologue : la fonction exponentielle
\item Lecture : Walter Rudin, \textit{Principles of Mathematical Analysis}, ch. 8
\end{refs}

\bigskip

\begin{problem}[{La sommation par parties, et le théorème d'Abel --- Master, short}]
\textbf{CAL-41.} \textbf{Problème.} Démontrez le théorème d'Abel : si $ \sum_{n\ge0} a_n $
converge, alors la série entière $ f(x) = \sum_{n\ge 0} a_n x^n $ converge pour
$ 0 \le x < 1 $ et
\[ \lim_{x\to 1^{-}} f(x) \;=\; \sum_{n=0}^{\infty} a_n . \]
Énoncez et démontrez d'abord l'identité de sommation par parties et servez-vous-en comme moteur ; déduisez-en ensuite les
sommes de la série harmonique alternée et de la série de Leibniz
$ 1 - \frac13 + \frac15 - \cdots $ à partir des séries entières de $ \log(1+x) $ et de
$ \arctan x $.
\end{problem}

\noindent Niels Henrik Abel a démontré ce résultat dans son mémoire de 1826 sur la série du binôme, l'année
même où il écrivait à Holmboe que les séries divergentes sont \og{} une invention du diable \fg{}
et qu'il est honteux de fonder une démonstration sur elles. Le théorème est le permis d'une manœuvre
exécutée constamment et d'ordinaire sans commentaire --- évaluer une série au bord de son disque de
convergence. La réciproque est fausse, et chercher des hypothèses sous lesquelles elle devient vraie a lancé
toute la théorie des théorèmes taubériens, à commencer par Tauber en 1897 et Littlewood en
1911.

\begin{refs}
\item Contexte : Théorème d'Abel --- Wikipédia (\url{https://fr.wikipedia.org/wiki/Th%C3%A9or%C3%A8me_d%27Abel_%28analyse%29})
\item Contexte : Sommation par parties --- Wikipédia (\url{https://fr.wikipedia.org/wiki/Sommation_par_parties})
\item Lecture : Walter Rudin, \textit{Principles of Mathematical Analysis}, ch. 8
\end{refs}

\bigskip


\section*{Recherche}

\begin{problem}[{La constante d'Euler est-elle irrationnelle ? --- Recherche}]
\textbf{CAL-22.} \textbf{Problème.} La constante d'Euler est
\[ \gamma = \lim_{n\to\infty}\Big( 1 + \frac12 + \cdots + \frac1n - \ln n \Big)
\approx 0{,}57721, \]
dont vous avez démontré l'existence en CAL-11. \textbf{Ouvert :} $\gamma$ est-elle irrationnelle ? Est-elle transcendante ?
Personne ne le sait --- on ne sait même pas que $\gamma$ $\notin$ $\mathbb{Q}$. Sous-problèmes accessibles : vérifiez à partir de la
définition la représentation intégrale
$ \gamma = \int_0^1 \big( \frac{1}{1-x} + \frac{1}{\ln x} \big)\,dx $ ; lisez ensuite ce que l'on
sait \textit{vraiment} : les calculs de fractions continues montrent que si $ \gamma = p/q $ alors
$ q $ est astronomiquement grand (plus de $ 10^{240000} $), et Rivoal a démontré (2012)
qu'au moins l'une des deux constantes $\gamma$ et la constante d'Euler--Gompertz
$ \delta = \int_0^\infty \frac{e^{-x}}{1+x}\,dx $ est transcendante --- sans décider
laquelle.
\end{problem}

\noindent La constante est celle d'Euler (1734) ; elle traverse la fonction Gamma,
la fonction zêta et la théorie analytique des nombres. Hilbert qualifiait son irrationalité
d'\og{} inabordable \fg{} ; Conway et Guy ont écrit qu'ils étaient prêts à parier qu'elle est
transcendante mais qu'ils n'espéraient pas de démonstration de leur vivant. Statut en 2026 : entièrement
ouvert, et parmi les problèmes ouverts les plus célèbres portant sur un nombre donné. Voir aussi
Problèmes ouverts (\url{open-problems.html}).

\begin{refs}
\item Contexte : Constante d'Euler-Mascheroni --- Wikipédia (\url{https://fr.wikipedia.org/wiki/Constante_d%27Euler-Mascheroni})
\item Lecture : Julian Havil, \textit{Gamma: Exploring Euler's Constant}
\item Article : J. C. Lagarias, \og{} Euler's constant: Euler's work and modern developments \fg{}, Bull. AMS 50 (2013), arXiv:1303.1856 (\url{https://arxiv.org/abs/1303.1856})
\end{refs}

\bigskip

\begin{problem}[{L'intégration en termes finis --- Recherche}]
\textbf{CAL-23.} \textbf{Problème.} Pourquoi $ \int e^{-x^2}dx $ n'a-t-elle pas de formule ? Rendez
la question précise et attaquez-la :
\begin{enumerate}
\item[(a)] définissez les fonctions élémentaires (le corps
engendré à partir des fractions rationnelles par les exponentielles, les logarithmes et les opérations algébriques) ;
\item[(b)] énoncez le théorème de Liouville sur la forme que doit prendre une primitive élémentaire ;
\item[(c)] servez-vous-en pour démontrer que $ e^{x^2} $ n'a pas de primitive élémentaire, en montrant que
l'équation différentielle requise $ R' + 2xR = 1 $ n'a pas de solution en fractions rationnelles
$ R $. Parcourez ensuite la frontière : l'algorithme de Risch (1969) décide de
l'intégrabilité élémentaire pour de larges classes d'intégrandes, mais une implémentation complète couvrant le
cas général mixte algébrique--transcendant n'a jamais été achevée --- le second volume que
Bronstein projetait n'était pas écrit à sa mort en 2005, et les systèmes de calcul formel
font encore tourner des versions partielles.
\end{enumerate}

\end{problem}

\noindent Liouville a démontré les premiers théorèmes d'impossibilité de ce type dans les années 1830 ---
l'analogue, pour l'intégration, de l'insolubilité de la quintique. Risch a transformé la question
en procédure de décision un siècle plus tard, un jalon précoce du calcul formel. Statut :
la théorie est décidable en principe pour les intégrandes élémentaires ; rendre l'algorithme pleinement
explicite et l'implémenter reste une frontière mathématique et technique active, non une
conjecture ouverte.

\begin{refs}
\item Contexte : Algorithme de Risch --- Wikipédia (\url{https://fr.wikipedia.org/wiki/Algorithme_de_Risch})
\item Contexte : Théorème de Liouville (algèbre différentielle) --- Wikipédia (\url{https://fr.wikipedia.org/wiki/Th%C3%A9or%C3%A8me_de_Liouville_%28alg%C3%A8bre_diff%C3%A9rentielle%29})
\item Contexte : Intégrale non élémentaire --- Wikipédia (\url{https://fr.wikipedia.org/wiki/Int%C3%A9grale_non_%C3%A9l%C3%A9mentaire})
\item Lecture : Manuel Bronstein, \textit{Symbolic Integration I} (Springer)
\end{refs}

\bigskip

\begin{problem}[{Les périodes : les nombres que fabrique l'intégration --- Recherche}]
\textbf{CAL-24.} \textbf{Problème.} Une \textit{période} est un nombre complexe dont les parties réelle et
imaginaire sont des valeurs d'intégrales absolument convergentes de fractions rationnelles à
coefficients rationnels, sur des domaines de $ \mathbb{R}^n $ découpés par des inégalités
polynomiales à coefficients rationnels.
\begin{enumerate}
\item[(a)] Montrez que $ \sqrt{2} $, $ \ln 2 $, $\pi$
et $ \zeta(2) $ sont des périodes --- pour $\pi$, intégrez $ 1 $ sur
$ \{x^2+y^2\le 1\} $.
\item[(b)] Montrez que les périodes forment un anneau dénombrable contenant tous les
nombres algébriques --- de sorte qu'une infinité non dénombrable de réels ne sont pas des périodes, mais démontrer qu'un
nombre \textit{donné} naturel n'est pas une période est une tout autre affaire.
\textbf{Ouvert :} la conjecture de Kontsevich--Zagier --- deux représentations intégrales quelconques d'une
même période sont reliées par une chaîne de mouvements élémentaires (additivité, changement de variables,
Newton--Leibniz) --- est ouverte ; on conjecture, sans l'avoir démontré, que $ e $ et $ 1/\pi $
ne sont pas des périodes ; et savoir si $\gamma$ (CAL-22) est une période est de même inconnu.
\end{enumerate}

\end{problem}

\noindent Kontsevich et Zagier ont introduit l'anneau des périodes dans leur essai de 2001
comme \og{} la classe de nombres la plus importante entre les nombres algébriques et tous les nombres
complexes \fg{} --- les nombres que l'analyse produit réellement. La conjecture rendrait l'égalité
de périodes vérifiable algorithmiquement en principe et engloberait les énoncés classiques de
transcendance ; statut en 2026 : grand ouvert, avec de profondes connexions aux motifs
en géométrie algébrique.

\begin{refs}
\item Contexte : Période (théorie des nombres) --- Wikipédia (\url{https://fr.wikipedia.org/wiki/Alg%C3%A8bre_des_p%C3%A9riodes})
\item Article : M. Kontsevich \& D. Zagier, \og{} Periods \fg{}, dans \textit{Mathematics Unlimited --- 2001 and Beyond}, Springer (2001)
\end{refs}

\bigskip

\begin{problem}[{$ \pi + e $ est-il irrationnel ? --- Recherche, short}]
\textbf{CAL-32.} \textbf{Problème.} On sait que $ \pi $ et $ e $ sont tous deux
transcendants, et pourtant on \textbf{ne sait pas} si $ \pi + e $ est irrationnel --- ni
si $ \pi e $, $ \pi/e $ ou $ \pi^{e} $ le sont. Démontrez la seule chose facile qui
\textit{puisse} se dire : au moins l'un des deux nombres $ \pi + e $ et $ \pi e $ est transcendant.
\end{problem}

\noindent La démonstration du fait énoncé tient en trois lignes --- considérez le polynôme
dont les racines sont $ \pi $ et $ e $, et utilisez que les nombres algébriques forment
un corps --- et le contraste avec les questions ouvertes est la leçon. Les démonstrations de transcendance sont
délicates et essentiellement uniques : celle de Lindemann pour $ \pi $ en 1882 et celle d'Hermite pour
$ e $ en 1873 sont des jalons, et plus d'un siècle plus tard la somme de ces deux nombres
nous met en échec. La conjecture de Schanuel réglerait tout cela d'un coup, ce qui donne une bonne
mesure de ce qui repose sur elle.

\begin{refs}
\item Contexte : Nombre transcendant --- Wikipédia (\url{https://fr.wikipedia.org/wiki/Nombre_transcendant})
\item Contexte : Conjecture de Schanuel --- Wikipédia (\url{https://fr.wikipedia.org/wiki/Conjecture_de_Schanuel})
\item Voir aussi : Problèmes ouverts (\url{open-problems.html}), Théorie des nombres (\url{pure-number-theory.html})
\item Voir aussi : OPEN-64 (\url{open-problems.html#OPEN-64}) --- la page Problèmes ouverts, où il figure parmi les questions classiques non résolues.
\end{refs}

\bigskip

\begin{problem}[{La série de Flint Hills --- Recherche, short}]
\textbf{CAL-33.} \textbf{Problème.} Déterminez si
\[ \sum_{n=1}^{\infty} \frac{1}{n^{3} \sin^{2} n} \]
converge. Statut : \textbf{ouvert}. Montrez que la convergence découlerait d'une majoration suffisamment
forte de la finesse avec laquelle $ \pi $ peut être approché par des rationnels --- c'est-à-dire de la
mesure d'irrationalité de $ \pi $ --- et énoncez la meilleure borne connue à ce jour.
\end{problem}

\noindent Chaque terme est fini, mais $ \sin n $ s'approche arbitrairement près de zéro
dès que $ n $ est proche d'un multiple de $ \pi $, de sorte que la série vit ou meurt selon la qualité avec
laquelle $ \pi $ peut être approché par des fractions. C'est ce qui fait d'une innocente
question de convergence, du genre de celles posées en première année d'analyse, l'équivalent d'un problème ouvert
difficile d'approximation diophantienne. C'est la démonstration la plus nette de ce site que
\og{} cette série converge-t-elle ? \fg{} n'est pas toujours un exercice de routine.

\begin{refs}
\item Contexte : Série de Flint Hills --- Wikipédia (en anglais) (\url{https://en.wikipedia.org/wiki/Flint_Hills_series})
\item Contexte : Mesure d'irrationalité --- Wikipédia (en anglais) (\url{https://en.wikipedia.org/wiki/Irrationality_measure})
\item Voir aussi : Analyse réelle (\url{pure-real-analysis.html}), Problèmes ouverts (\url{open-problems.html})
\end{refs}

\bigskip

\begin{problem}[{Une intégrale qui intègre toute dérivée --- Recherche}]
\textbf{CAL-42.} \textbf{Problème.} CAL-09(iii) produit une fonction dérivable en tout point de $ [0,1] $
dont la dérivée n'est pas intégrable au sens de Riemann, et l'intégrale de Lebesgue ne répare pas cela :
$ F(x) = x^2\sin(1/x^2) $, $ F(0)=0 $, est partout dérivable avec
$ \int_0^1 |F'| = \infty $. Une intégrale, elle, répare la chose. Étant donné une fonction strictement positive
$ \delta $ sur $ [a,b] $ --- une \textit{jauge} --- on dit qu'une subdivision pointée
$ a = x_0 < \cdots < x_n = b $ de points marqués $ t_i \in [x_{i-1},x_i] $ est \textit{$ \delta $-fine}
si $ [x_{i-1},x_i] \subset (t_i-\delta(t_i),\, t_i+\delta(t_i)) $ pour tout $ i $ ; alors $ f $
est intégrable au sens de Henstock--Kurzweil d'intégrale $ I $ si pour tout $ \varepsilon > 0 $
il existe une jauge $ \delta $ telle que toute subdivision pointée $ \delta $-fine ait une somme de Riemann
à moins de $ \varepsilon $ de $ I $.
\begin{enumerate}
\item[(a)] Démontrez le lemme de Cousin : pour toute jauge $ \delta $ sur $ [a,b] $ il existe une subdivision
pointée $ \delta $-fine. (Dichotomisez, et utilisez la complétude.) Sans cela la définition serait
vide.
\item[(b)] Démontrez le théorème fondamental sous sa forme parfaite : si $ F $ est dérivable en tout
point de $ [a,b] $, alors $ F' $ est intégrable au sens de Henstock--Kurzweil et
$ \int_a^b F' = F(b)-F(a) $ --- sans aucune hypothèse d'intégrabilité sur $ F' $.
\item[(c)] Démontrez que toute fonction intégrable au sens de Lebesgue est intégrable au sens de Henstock--Kurzweil avec la
même valeur, et que $ f $ est intégrable au sens de Lebesgue si et seulement si $ f $ et $ |f| $ sont toutes deux
intégrables au sens de Henstock--Kurzweil. Concluez que cette intégrale est strictement plus générale et n'est
pas une intégrale absolue, et exhibez une fonction qu'elle n'intègre que de façon conditionnelle.
\item[(d)] La frontière. Formulez le problème en dimension $ n $ : trouver une intégrale pour laquelle le
théorème de la divergence $ \int_{\Omega} \mathrm{div}\, v = \int_{\partial\Omega} v\cdot n $ vaille pour
tout champ de vecteurs ponctuellement dérivable sur tout ensemble borné de périmètre fini. Lisez
l'exposé de Pfeffer sur ce qui a été accompli, et rédigez une appréciation des raisons pour lesquelles l'intégrale
de jauge n'a jamais supplanté celle de Lebesgue en pratique --- ce qu'il advient des théorèmes de
convergence, de Fubini, et de la mesure sous-jacente.
\end{enumerate}

\end{problem}

\noindent La théorie en dimension un est réglée et magnifique : Arnaud Denjoy a construit en 1912, par
récursion transfinie, une intégrale inversant toute dérivée, Perron en a donné une équivalente
en 1914 au moyen de fonctions majorantes et minorantes, et toutes deux ont été jugées d'une technicité rebutante jusqu'à ce que
Jaroslav Kurzweil (1957) et Ralph Henstock (1961) découvrent que remplacer le $ \delta $ constant
de la définition de Riemann par une fonction du point marqué donne exactement la même classe
de fonctions intégrables, avec une définition qu'un étudiant de première année peut lire. \textbf{Statut :} les parties
(a)--(c) sont des théorèmes ; la partie (d) est un véritable domaine de recherche --- les intégrales descriptives et
de jauge multidimensionnelles de Jarník, Kurzweil, Mařík, Pfeffer et De Pauw démontrent bien des théorèmes de la
divergence très généraux, mais la théorie est technique, non canonique et encore en développement, et
la question pédagogique de savoir si l'intégrale de jauge devrait remplacer celle de Lebesgue dans l'enseignement
reste un débat vivant plutôt qu'une affaire réglée.

\begin{refs}
\item Contexte : Intégrale de Kurzweil--Henstock --- Wikipédia (\url{https://fr.wikipedia.org/wiki/Int%C3%A9grale_de_Kurzweil-Henstock})
\item Contexte : Théorème de la divergence --- Wikipédia (\url{https://fr.wikipedia.org/wiki/Th%C3%A9or%C3%A8me_de_la_divergence})
\item Lecture : Robert G. Bartle, \textit{A Modern Theory of Integration}, AMS Graduate Studies in Mathematics 32 (2001)
\item Lecture : Washek F. Pfeffer, \textit{Derivation and Integration}, Cambridge Tracts in Mathematics (2001)
\end{refs}

\bigskip

\begin{problem}[{Reconnaître zéro --- Recherche, short}]
\textbf{CAL-43.} \textbf{Problème.} Richardson a démontré en 1968 qu'il n'existe aucun algorithme qui,
étant donné une expression construite à partir des nombres rationnels, de $ \pi $, de $ \ln 2 $ et d'une variable
$ x $ à l'aide de $ +,-,\times $, de la composition et des fonctions $ \sin $, $ \exp $ et
$ |\cdot| $, décide si cette expression s'annule identiquement sur $ \mathbb{R} $ --- lisez la
démonstration et sachez esquisser comment un problème diophantien y est encodé. Attaquez ensuite le
problème des \textit{constantes}, qui reste ouvert : énoncez précisément ce que serait un algorithme de reconnaissance de zéro
pour les constantes élémentaires, démontrez ce que l'argument de Richardson de 1997 vous donne si l'on
admet la conjecture de Schanuel, et déterminez ce que l'on sait inconditionnellement.
\end{problem}

\noindent C'est l'ombre portée de CAL-23 : l'algorithme de Risch pour l'intégration symbolique
a besoin, comme sous-programme, de savoir décider si une constante qu'il a construite est nulle, de sorte que
son statut d'algorithme dépend exactement de cette question. Décider la nullité paraît la chose la plus
élémentaire que l'on puisse demander d'une formule, et en présence de $ \sin $ c'est prouvablement
impossible ; ôtez les fonctions trigonométriques et le problème devient conjecturalement
décidable, reposant sur la conjecture de Schanuel --- la même conjecture de transcendance qui
réglerait la question de l'irrationalité de $ \pi + e $ (CAL-32). \textbf{Statut :} indécidable pour
la classe de Richardson ; ouvert, et conditionnellement décidable, pour les constantes élémentaires.

\begin{refs}
\item Contexte : Théorème de Richardson --- Wikipédia (\url{https://fr.wikipedia.org/wiki/Th%C3%A9or%C3%A8me_de_Richardson})
\item Contexte : Problème des constantes --- Wikipédia (en anglais) (\url{https://en.wikipedia.org/wiki/Constant_problem})
\item Article : D. Richardson, \og{} How to recognize zero \fg{}, \textit{Journal of Symbolic Computation} 24 (1997), 627--645
\item Article : D. Richardson, \og{} Some undecidable problems involving elementary functions of a real variable \fg{}, \textit{Journal of Symbolic Logic} 33 (1968), 514--520
\end{refs}

\bigskip


\section*{Pour aller plus loin}


\vfill
\noindent\emph{Hors de la Caverne} --- des probl\`emes, pas de r\'eponses.
\end{document}
